\documentclass{article}
\usepackage[T1]{fontenc}
\usepackage{lmodern}
\usepackage{graphicx}
\usepackage{amsmath,amssymb}
\usepackage[a4paper,margin=2cm]{geometry}
\usepackage[absolute,overlay]{textpos}
\setlength{\TPHorizModule}{1mm}
\setlength{\TPVertModule}{1mm}
\usepackage[indonesian]{babel}
\usepackage{hyperref}
\setlength{\emergencystretch}{2em}
% unit: Halaman 94

\begin{document}

% === Halaman 94 (llm) ===

\begin{itemize}
    \item Nilai terbesar yang dapat muncul untuk merepresentasikan blok: $25252525$ (ZZZZ), maka $25252525$ dapat digunakan sebagai modulus $n$.
    \item Nilai $m$ yang relatif prima dengan $25252525$, misalnya $21035433$,
    \item $b$ dipilih antara $1$ dan $25252525$, misalnya $23210025$.
\end{itemize}

\noindent Fungsi enkripsi menjadi:
\[ C \equiv 21035433P + 23210025 \pmod{25252525} \]

\noindent Fungsi dekripsi, setelah dihitung, menjadi
\[ P \equiv 5174971 (C - 23210025) \pmod{25252525} \]

\end{document}
